\section*{अध्याय \ref{ch:b1:vspaces} --- सदिश समष्टियाँ}

\begin{solution}{exo:b1:vspaces:1}
\begin{enumerate}
  \item हाँ: उसमें $0$ है, और परिभाषक समीकरण रैखिक है ($x + \lambda y$ के
  अंतर्गत स्थायी)।
  \item नहीं: उसमें $(0,0,0)$ नहीं है।
  \item नहीं: $(1, 0)$ और $(0, -1)$ उसमें हैं ($xy = 0$), पर उनका योग $(1,
  -1)$ नहीं ($xy = -1 < 0$)।
  \item हाँ: $0$ $1$ पर शून्य होता है; $(P + \lambda Q)(1) = P(1) + \lambda
  Q(1) = 0$।
  \item हाँ: शून्य फलन परिबद्ध है; और यदि $\abs f \leq M$ तथा $\abs g \leq
  M'$, तो $\abs{f + \lambda g} \leq M + \abs\lambda M'$।
\end{enumerate}
\end{solution}

\begin{solution}{exo:b1:vspaces:2}
$(1,2,1) = a(1,0,1) + b(1,1,0)$ के लिए $a + b = 1$, $b = 2$, $a = 1$ चाहिए:
जो असंगत है ($a + b = 3 \neq 1$): अतः \omterm{def:b1:vspaces:span}{जनित समष्टि} में नहीं। $(2,1,1) =
a(1,0,1) + b(1,1,0)$: $b = 1$, $a = 1$, $a + b = 2$: संगत, अतः $(2,1,1) =
(1,0,1) + (1,1,0)$, और वह \omterm{def:b1:vspaces:span}{जनित समष्टि} में है।

समीकरण: $(x, y, z) = (a + b, b, a)$ का अर्थ है $x = y + z$: अतः \omterm{def:b1:vspaces:span}{जनित समष्टि}
समतल $\{x - y - z = 0\}$ है।
\end{solution}

\begin{solution}{exo:b1:vspaces:3}
पहला कुल: $\lambda(1,1,0) + \mu(1,0,1) + \nu(0,1,1) = 0$ से $\lambda + \mu =
0$, $\lambda + \nu = 0$, $\mu + \nu = 0$ मिलता है: जोड़ने पर $2(\lambda +
\mu + \nu) = 0$, और हर समीकरण को घटाने पर $\lambda = \mu = \nu = 0$: अतः
\omterm{def:b1:vspaces:free}{स्वतंत्र}।

दूसरा: $(2,4,6) = 2(1,2,3)$: परतंत्र।

तीसरा: $\R^3$ में चार सदिश --- विमा उपलब्ध होते ही अनिवार्यतः परतंत्र
(\cref{ch:b1:findim}); सीधे: $(0,1,1) = -(1,0,0) + 0\cdot(1,1,0) + (1,1,1)$,
वस्तुतः $(-1,0,0) + (1,1,1) = (0,1,1)$: अर्थात् एक अतुच्छ संबंध।
\end{solution}

\begin{solution}{exo:b1:vspaces:4}
बहुपदों $1, (X-1), (X-1)^2, (X-1)^3$ की घातें भिन्न-भिन्न हैं $0, 1, 2, 3$:
अतः \omterm{def:b1:vspaces:free}{स्वतंत्र} (\cref{prop:b1:vspaces:freecriteria}~(1)); और चार \omterm{def:b1:vspaces:free}{स्वतंत्र} सदिश
जनक भी हैं (हर $P \in \R_3[X]$ $X - 1$ की घातों में प्रसारित हो जाता है,
जैसे बहुपदों के लिए टेलर से; \cref{prop:b1:poly:multiplicity} की उपपत्ति
देखिए): अतः \omterm{def:b1:vspaces:free}{आधार}। $X^3$ के लिए $1$ पर टेलर: $P = X^3$, $P(1) = 1$, $P'(1) =
3$, $P''(1) = 6$, $P'''(1) = 6$:
\[
X^3 = 1 + 3(X - 1) + 3(X-1)^2 + (X-1)^3 ,
\]
\omterm{prop:b1:vspaces:coordinates}{निर्देशांक} $(1, 3, 3, 1)$ (पास्कल की पंक्ति, जैसी $X^3 = ((X-1)+1)^3$ से
अपेक्षित थी)।
\end{solution}

\begin{solution}{exo:b1:vspaces:5}
$F \cap G$: प्रतिबंध $x = y = z$ और $x = t = 0$ मिलकर $x = 0$ देते हैं, अतः
$y = z = 0$, और $t = 0$: सर्वनिष्ठ $\{0\}$ है। योग: $(x,y,z,t)$ दिया हो, तो
$(a,a,a,b) \in F$ और $(0,c,d,0) \in G$ ऐसे खोजिए जिनका योग वही हो: $a = x$,
$b = t$, $c = y - x$, $d = z - x$: यह सदा संभव है। अतः $\R^4 = F \oplus G$,
और
\[
(1,2,3,4) = (1,1,1,4) + (0,1,2,0) .
\]
\end{solution}

\begin{solution}{exo:b1:vspaces:6}
$G$ का कोई अवयव $\lambda u$ है; यदि वह अभिसरित होता है, तो (चूँकि $\lambda
u_n = \lambda(-1)^n$ के दो उपानुक्रमीय सीमा-मान $\pm \lambda$ हैं)
अनिवार्यतः $\lambda = 0$: $F \cap G = \{0\}$।

$F + G$ सब कुछ नहीं है: उसमें $(c_n)$ के अभिसारी होने पर $c_n +
\lambda(-1)^n$ रूप के अनुक्रम हैं। अनुक्रम $v_n = n$ इस रूप का नहीं है ($v_n
- \lambda(-1)^n$ अपरिबद्ध है, अतः कभी अभिसारी नहीं)। अतः $F \oplus G
\subsetneq$ (सभी अनुक्रमों की समष्टि)।
\end{solution}

\begin{solution}{exo:b1:vspaces:7}
($\supseteq$) $F \cap G$ और $F \cap H$ दोनों $F$ में हैं, और उनका योग $G +
(F \cap H)$ में है: अंतर्विष्टता इसी से आ जाती है, क्योंकि बायाँ पक्ष दोनों
टुकड़ों को समाहित करने वाली \omterm{def:b1:vspaces:subspace}{उपसमष्टि} है --- ठोस रूप से, $g \in F\cap G$, $h
\in F \cap H$ के साथ कोई अवयव $g + h$ $F$ में है ($F$ के दो अवयवों का योग)
और $G + (F \cap H)$ में भी।

($\subseteq$) मान लीजिए $x = g + h$, $g \in G$, $h \in F \cap H$ के साथ $x
\in F$। तब $g = x - h \in F$ ($F$ के अवयवों का अंतर), अतः $g \in F \cap G$,
और $x = g + h \in (F \cap G) + (F \cap H)$।

$\R^2$ में पूर्ण वितरण नियम का प्रति-उदाहरण: $F = \operatorname{Vect}(1,1)$,
$G = \operatorname{Vect}(1,0)$, $H = \operatorname{Vect}(0,1)$। तब $G + H =
\R^2$, अतः $F \cap (G+H) = F$, जबकि $F \cap G = F \cap H = \{0\}$: दायाँ
पक्ष $\{0\} \neq F$ है।
\end{solution}

\begin{solution}{exo:b1:vspaces:8}
\begin{enumerate}
  \item मान लीजिए सभी $x$ के लिए $\sum_{i} \lambda_i \eu^{a_i x} = 0$।
  $\eu^{-a_p x}$ से गुणा कीजिए: $x \to +\infty$ होने पर $\lambda_p + \sum_{i
  < p} \lambda_i \eu^{(a_i - a_p)x} \to \lambda_p$ (हर घातांक $a_i - a_p <
  0$)। बायाँ पक्ष सर्वसम रूप से $0$ है, अतः $\lambda_p = 0$; और यही नीचे की
  ओर दोहराइए।
  \item मान लीजिए सभी $x$ के लिए $a\cos x + b \sin x + c \cos 2x + d\sin 2x
  = 0$। $x = 0$ पर मूल्यांकन कीजिए: $a + c = 0$; $x = \pi$ पर: $-a + c = 0$;
  अतः $a = c = 0$, और संबंध घटकर $b\sin x + d \sin 2x = 0$ रह जाता है। $x =
  \frac\pi2$ पर मूल्यांकन कीजिए: $b = 0$; फिर $x = \frac\pi4$ पर: $d = 0$।
  \item मान लीजिए सभी $x$ के लिए $\sum \lambda_i \abs{x - a_i} = 0$। फलन
  $\sum_{i \neq j} \lambda_i\abs{x - a_i}$ $a_j$ पर \omterm{def:b1:derivative:def}{अवकलनीय} है (हर पद अपने
  कोने से दूर \omterm{def:b1:derivative:def}{अवकलनीय} है), अतः उनका अंतर $-\lambda_j \abs{x - a_j}$ भी $a_j$
  पर \omterm{def:b1:derivative:def}{अवकलनीय} होना ही चाहिए --- और यह $\lambda_j = 0$ को बाध्य कर देता है
  ($\abs{\,\cdot\,}$ में एक कोना है)। यह हर $j$ के लिए लागू होता है।
\end{enumerate}
\end{solution}

\begin{solution}{exo:b1:vspaces:9}
मान लीजिए $h \in H$। चूँकि $h \in H \subseteq F + H = F + G$, अतः $f \in F$,
$g \in G$ के साथ $h = f + g$ लिखिए। तब $f = h - g \in H$ (दोनों पद $H$ में
हैं, $G \subseteq H$ का प्रयोग करते हुए), अतः $f \in F \cap H = F \cap G
\subseteq G$, और $h = f + g \in G$। अतः $H \subseteq G$, और परिकल्पना $G
\subseteq H$ के साथ: समानता।

$G \subseteq H$ के बिना प्रति-उदाहरण: $\R^2$ में $F =
\operatorname{Vect}(1,0)$, $G = \operatorname{Vect}(0,1)$, $H =
\operatorname{Vect}(1,1)$ लीजिए: तब $F + G = F + H = \R^2$ और $F \cap G = F
\cap H = \{0\}$, फिर भी $G \neq H$।
\end{solution}

\begin{solution}{exo:b1:vspaces:10}
दोनों \omterm{def:b1:logic:sets}{समुच्चय} उपसमष्टियाँ हैं (परिभाषक प्रतिबंध रैखिक हैं और $0$ के लिए सत्य
हैं)। विघटन: $P \in \R[X]$ के लिए,
\[
P(X) = \underbrace{\frac{P(X) + P(-X)}{2}}_{\in\,\mathcal P}
+ \underbrace{\frac{P(X) - P(-X)}{2}}_{\in\,\mathcal I},
\]
और जो \omterm{def:b1:poly:def}{बहुपद} सम भी हो और विषम भी, वह $P = -P$ संतुष्ट करता है, अतः $P = 0$:
इस प्रकार योग प्रत्यक्ष है और $\R[X]$ के बराबर है।

अब मान लीजिए $P = \sum_k a_k X^k$ सम है। तब $P(X) - P(-X) = 2 \sum_{k \text{
विषम}} a_k X^k$ शून्य \omterm{def:b1:poly:def}{बहुपद} है, अतः विषम घात वाला हर गुणांक शून्य हो जाता है
(\cref{def:b1:poly:def}): $P \in \operatorname{Vect}(1, X^2, X^4, \dots)$,
अर्थात् किसी \omterm{def:b1:poly:def}{बहुपद} $Q$ के लिए $P = Q(X^2)$। विलोमतः $X^2$ का हर \omterm{def:b1:poly:def}{बहुपद} सम है।
\end{solution}

\begin{solution}{exo:b1:vspaces:11}
मान लीजिए $a(x_1 + x_2) + b(x_2 + x_3) + c(x_3 + x_1) = 0$। \omterm{def:b1:vspaces:free}{स्वतंत्र कुल}
$(x_1, x_2, x_3)$ पर फिर से समूहबद्ध करने पर:
\[
(a + c)\,x_1 + (a + b)\,x_2 + (b + c)\,x_3 = 0
\implies a + c = a + b = b + c = 0 .
\]
पहले दो समीकरणों को घटाने पर $c = b$ मिलता है; फिर तीसरे से $2b = 0$, अतः $b
= c = 0$, और फिर $a = 0$: अतः कुल \omterm{def:b1:vspaces:free}{स्वतंत्र} है।

चार सदिशों के लिए तदनुरूप कुल \emph{सदा} परतंत्र होता है:
\[
(x_1 + x_2) - (x_2 + x_3) + (x_3 + x_4) - (x_4 + x_1) = 0
\]
एक अतुच्छ शून्य संयोजन है (गुणांक $1, -1, 1, -1$), $(x_1, x_2, x_3, x_4)$
चाहे जो भी हो। निर्णय चक्र की लंबाई की समता करती है।
\end{solution}

\begin{solution}{exo:b1:vspaces:12}
\begin{enumerate}
  \item यदि $F_1 \subseteq F_2$ अथवा $F_2 \subseteq F_1$, तो संघ दोनों में
  से एक ही है, अतः उचित है। अन्यथा $x \in F_1 \setminus F_2$ और $y \in F_2
  \setminus F_1$ चुनिए, और $x + y$ पर विचार कीजिए। यदि $x + y \in F_1$, तो
  $y = (x + y) - x \in F_1$: विरोधाभास। यदि $x + y \in F_2$, तो $x \in F_2$:
  विरोधाभास। अतः $x + y \notin F_1 \cup F_2$, और $E \neq F_1 \cup F_2$।
  \item विरोधाभास के लिए मान लीजिए $E = F_1 \cup \dots \cup F_k$, जहाँ $k$
  ऐसे सभी आच्छादनों में \emph{न्यूनतम} चुना गया है। न्यूनतमता $F_1 \subseteq
  F_2 \cup \dots \cup F_k$ को मना कर देती है (अन्यथा $F_1$ हटा दीजिए), अतः
  ऐसा $x \in F_1$ है कि सभी $i \geq 2$ के लिए $x \notin F_i$। चूँकि $F_1$
  उचित है, $y \notin F_1$ चुनिए। हर अदिश $t$ के लिए सदिश $y + t x$ किसी न
  किसी $F_i$ में है। वह कभी $F_1$ में नहीं होता: अन्यथा $y = (y + tx) - tx
  \in F_1$ (क्योंकि $x \in F_1$)। \omterm{def:b1:structures:field}{क्षेत्र} अनंत है, अतः $k$ भिन्न-भिन्न अदिश
  $t_1, \dots, t_k$ चुनिए: $k$ सदिश $y + t_j x$ $k - 1$ उपसमष्टियों $F_2,
  \dots, F_k$ में गिरते हैं, और उनमें से दो, मान लीजिए $t \neq t'$ के साथ $y
  + t x$ और $y + t' x$, एक ही $F_i$ में होंगे ($i \geq 2$)। तब उनका अंतर $(t
  - t')x \in F_i$, अतः $x \in F_i$: विरोधाभास। अतः उचित उपसमष्टियों द्वारा
  कोई परिमित आच्छादन है ही नहीं।
\end{enumerate}
\end{solution}

\begin{solution}{pb:b1:vspaces:1}
\textbf{1.} $L_i$ $n$ रैखिक गुणनखंडों का गुणनफल है, जिसे एक अशून्य अचर से
भाग दिया गया है ($x_i$ भिन्न-भिन्न हैं), अतः $\deg L_i = n$। $j \neq i$ के
साथ $x_j$ पर मूल्यांकन करने पर: अंश का गुणनखंड $X - x_j$ शून्य हो जाता है,
अतः $L_i(x_j) = 0$। $x_i$ पर अंश और हर एक ही हैं: $L_i(x_i) = 1$।

\textbf{2.} मान लीजिए $\sum_i \lambda_i L_i = 0$। $x_j$ पर मूल्यांकन कीजिए:
$\lambda_j L_j(x_j) = \lambda_j$ को छोड़कर सभी पद मर जाते हैं, अतः हर $j$ के
लिए $\lambda_j = 0$: कुल \omterm{def:b1:vspaces:free}{स्वतंत्र} है।

\textbf{3.} मान लीजिए $D = P - \sum_i P(x_i) L_i$। तब $\deg D \leq n$ और
प्रश्न 1 से $n + 1$ भिन्न-भिन्न बिंदुओं $x_0, \dots, x_n$ के लिए $D(x_j) =
P(x_j) - P(x_j) = 0$। $\leq n$ घात वाले किसी अशून्य \omterm{def:b1:poly:def}{बहुपद} के अधिक से अधिक
$n$ मूल होते हैं (\cref{cor:b1:poly:nroots}), अतः $D = 0$। इस प्रकार हर $P
\in \R_n[X]$ $L_i$ का संयोजन है: कुल जनक है, और प्रश्न 2 के साथ \omterm{def:b1:vspaces:free}{आधार}।

\textbf{4.} $y_0, \dots, y_n$ दिया हो, तो \omterm{def:b1:poly:def}{बहुपद} $P = \sum_i y_i L_i$ की घात
$\leq n$ है और वह अंतर्वेशन करता है। अद्वितीयता: प्रश्न 3 से किसी भी
अंतर्वेशी $P$ के \omterm{def:b1:vspaces:free}{आधार} $(L_i)$ में \omterm{prop:b1:vspaces:coordinates}{निर्देशांक} $(P(x_0), \dots, P(x_n)) =
(y_0, \dots, y_n)$ हैं, और किसी \omterm{def:b1:vspaces:free}{आधार} में \omterm{prop:b1:vspaces:coordinates}{निर्देशांक} अद्वितीय होते हैं
(\cref{prop:b1:vspaces:coordinates})। लाग्रांज \omterm{def:b1:vspaces:free}{आधार} में $P$ के \omterm{prop:b1:vspaces:coordinates}{निर्देशांक}
\emph{गाँठों पर उसके मान} ही हैं --- इस \omterm{def:b1:vspaces:free}{आधार} का पूरा मर्म यही है।

\textbf{5.} प्रश्न 3 को $P = X^k$ पर लगाइए ($0 \leq k \leq n$):
\[
X^k = \sum_{i=0}^{n} x_i^{k} L_i ,
\]
और $k = 0$ से $\sum_i L_i = 1$ मिलता है।

\textbf{6.} ठीक-ठीक $\deg N_k = k$: कुल $(N_0, \dots, N_n)$ $\R_n[X]$ में
घातों की एक सीढ़ी है, अतः \cref{ex:b1:vspaces:staircase} से \omterm{def:b1:vspaces:free}{आधार} है
(स्वतंत्रता \cref{prop:b1:vspaces:freecriteria}~(1) से, और जनक होना घात पर
परिमित अवरोहण से)।

\textbf{7.} पुनरावृत्ति से,
\[
f[x_0, x_1] = \frac{f(x_1) - f(x_0)}{x_1 - x_0},
\qquad
f[x_0, x_1, x_2] = \frac{f[x_1, x_2] - f[x_0, x_1]}{x_2 - x_0}.
\]
$f(x) = x^2$ के लिए:
\[
f[x_0, x_1] = \frac{x_1^2 - x_0^2}{x_1 - x_0} = x_0 + x_1,
\]
और फिर
\[
f[x_0, x_1, x_2]
= \frac{(x_1 + x_2) - (x_0 + x_1)}{x_2 - x_0}
= \frac{x_2 - x_0}{x_2 - x_0} = 1 .
\]

\textbf{8.} $\deg S \leq n$, क्योंकि $Q, R$ की घात $\leq n - 1$ है। $x_0$
पर: $S(x_0) = \frac{-(x_0 - x_n) R(x_0)}{x_n - x_0} = R(x_0) = f(x_0)$।
$x_n$ पर: $S(x_n) = \frac{(x_n - x_0) Q(x_n)}{x_n - x_0} = Q(x_n) = f(x_n)$।
किसी आंतरिक गाँठ $x_i$ ($1 \leq i \leq n-1$) पर $Q$ और $R$ दोनों मान
$f(x_i)$ लेते हैं, अतः
\[
S(x_i) = \frac{(x_i - x_0) - (x_i - x_n)}{x_n - x_0}\, f(x_i)
= f(x_i) .
\]

\textbf{9.} गाँठों की संख्या पर आगमन। एक गाँठ: अंतर्वेशक अचर $f(x_0) =
f[x_0]$ है। दावे को $k$ गाँठों के लिए मान लीजिए और $S$ को $x_0, \dots, x_k$
पर अंतर्वेशक लीजिए; अद्वितीयता (प्रश्न 4) से $S$ ऐटकिन की प्रमेयिका द्वारा
$R$ (गाँठें $x_0, \dots, x_{k-1}$) और $Q$ (गाँठें $x_1, \dots, x_k$) से
मिलता है। $S$ में $X^{k}$ का गुणांक है
\[
\frac{[X^{k-1}]\,Q - [X^{k-1}]\,R}{x_k - x_0}
= \frac{f[x_1, \dots, x_k] - f[x_0, \dots, x_{k-1}]}{x_k - x_0}
= f[x_0, \dots, x_k]
\]
--- आगमन-परिकल्पना और परिभाषक पुनरावृत्ति से।

\textbf{10.} मान लीजिए $P_k$ $x_0, \dots, x_k$ पर $f$ का अंतर्वेशन करता है।
अंतर $P_k - P_{k-1}$ की घात $\leq k$ है और वह $x_0, \dots, x_{k-1}$ पर शून्य
होता है, अतः गुणनखंड प्रमेय को $k$ बार लगाने पर (\cref{thm:b1:poly:factor})
वह किसी अचर $c$ के लिए $c\,N_k$ के बराबर है; $X^{k}$ के गुणांकों की तुलना
करने पर और प्रश्न 9 का प्रयोग करने पर $c = f[x_0, \dots, x_k]$। $P_0 =
f(x_0) N_0$ से दूरबीनी करने पर न्यूटन का सूत्र मिल जाता है। संवृत रूप के लिए
$P_k = \sum_{i \leq k} f(x_i) L_i$ लिखिए (लाग्रांज, गाँठों $x_0, \dots, x_k$
पर) और $X^{k}$ का गुणांक पढ़िए: हर $L_i$ $\frac{1}{\prod_{j \neq i}(x_i -
x_j)}$ का योगदान देता है, जिससे
\[
f[x_0, \dots, x_k] = \sum_{i=0}^{k}
\frac{f(x_i)}{\prod_{j \neq i,\, j \leq k}(x_i - x_j)} .
\]
दायाँ पक्ष गाँठों के किसी भी \omterm{def:b1:counting:objects}{क्रमचय} के अंतर्गत अपरिवर्तित है, अतः \omterm{pb:b1:vspaces:1}{विभाजित
अंतर} उनके क्रम पर निर्भर नहीं करता।

\textbf{11.} यदि $P = a X^m + (\text{निम्नतर घातें})$, तो द्विपद प्रमेय देती
है
\[
\Delta P = a\bigl((X+1)^m - X^m\bigr) + \dots
= a\,m\,X^{m-1} + (\text{निम्नतर घातें}),
\]
क्योंकि $(X+1)^m - X^m = m X^{m-1} + \dots$, और $P$ का निम्नतर-घात वाला भाग
$\Delta$ के बाद $\leq m - 2$ घात का योगदान देता है (अथवा $\leq m-2$ घात के
पद)। अतः $\deg \Delta P = m - 1$, जिसका अग्र गुणांक $m a$ है। किसी अचर $c$
के लिए $\Delta c = c - c = 0$।

\textbf{12.} $\deg B_k = k$: सीढ़ी, अतः $\R_n[X]$ का \omterm{def:b1:vspaces:free}{आधार}
(\cref{ex:b1:vspaces:staircase})। $\Delta B_k$ के लिए ($k \geq 1$) उभयनिष्ठ
गुणनफल का गुणनखंडन कीजिए:
\begin{align*}
k!\,\Delta B_k
&= (X+1)X\cdots(X-k+2) - X(X-1)\cdots(X-k+1) \\
&= X(X-1)\cdots(X-k+2)\,\bigl[(X+1) - (X-k+1)\bigr] \\
&= k\,X(X-1)\cdots(X-k+2),
\end{align*}
अतः $\Delta B_k = \frac{X(X-1)\cdots(X-k+2)}{(k-1)!} = B_{k-1}$।

\textbf{13.} $P = \sum_{k=0}^{n} c_k B_k$ लिखिए (\omterm{def:b1:vspaces:free}{आधार}, प्रश्न 12)।
$\Delta^{j}$ लगाइए: प्रश्न 12 से $\Delta^{j} P = \sum_{k \geq j} c_k
B_{k-j}$। $0$ पर मूल्यांकन कीजिए: $B_0(0) = 1$, और $m \geq 1$ के लिए $B_m(0)
= 0$ (गुणनखंड $X$ शून्य हो जाता है), अतः $\bigl(\Delta^{j}P\bigr)(0) = c_j$।
यही अग्र-अंतर सूत्र है।

\textbf{14.} $k$ पर आगमन। $k = 0$ के लिए सर्वसमिका $P(0) = P(0)$ कहती है।
उसे $k$ के लिए मान लीजिए और $\Delta P$ पर लगाइए:
\[
\bigl(\Delta^{k+1} P\bigr)(0)
= \sum_{j=0}^{k} (-1)^{k-j}\binom kj \bigl(P(j+1) - P(j)\bigr).
\]
$P(i)$ का गुणांक इकट्ठा कीजिए: पहले योग से (सरकाकर) वह $(-1)^{k-i+1}\binom
k{i-1}\cdot(-1)^{0}$ है और दूसरे से $-(-1)^{k-i}\binom ki$ --- दोनों मिलकर
\[
(-1)^{k+1-i}\Bigl(\binom k{i-1} + \binom ki\Bigr)
= (-1)^{k+1-i}\binom{k+1}i
\]
--- पास्कल के नियम से, जो कोटि $k + 1$ पर वही सर्वसमिका है।

\textbf{15.} प्रश्न 11 को घात $n$ और अग्र गुणांक $a_n$ से दोहराइए: एक
$\Delta$ के बाद घात $n-1$ और अग्र गुणांक $n a_n$; दो के बाद $n(n-1)a_n$; $n$
चरणों के बाद घात $0$ और मान $n(n-1)\cdots 1\, a_n = n!\,a_n$, जो एक अचर है।
एक और $\Delta$ उसे मार देता है: $\Delta^{n+1}P = 0$।

\textbf{16.} यदि $m \geq k$: $B_k(m) = \binom mk \in \N$। यदि $0 \leq m <
k$: $m(m-1)\cdots(m-k+1)$ का कोई एक गुणनखंड शून्य है, अतः $B_k(m) = 0$। यदि
$q \geq 1$ के साथ $m = -q$:
\[
B_k(-q) = \frac{(-q)(-q-1)\cdots(-q-k+1)}{k!}
= (-1)^k\,\frac{q(q+1)\cdots(q+k-1)}{k!}
= (-1)^k \binom{q+k-1}{k},
\]
जो एक पूर्णांक है। अतः हर $B_k$ $\Z$ को $\Z$ में भेजता है।

\textbf{17.} ($\Leftarrow$) यदि $c_k \in \Z$ के साथ $P = \sum_k c_k B_k$, तो
$m \in \Z$ के लिए प्रश्न 16 से $P(m) = \sum_k c_k B_k(m) \in \Z$।
($\Rightarrow$) यदि $P$ पूर्णांक-मान है, तो उसके \omterm{prop:b1:vspaces:coordinates}{निर्देशांक} $c_k =
\bigl(\Delta^k P\bigr)(0) = \sum_{j=0}^k (-1)^{k-j}\binom kj P(j)$ हैं
(प्रश्न 13 और 14), जो पूर्णांकों $P(0), \dots, P(k)$ का पूर्णांक संयोजन है।
यही पूर्णांक-मान बहुपदों का पोल्या द्वारा दिया गया अभिलक्षण है।

\textbf{18.} $Q(X) = P(X + a)$ रखिए, जो $\leq n$ घात का \omterm{def:b1:poly:def}{बहुपद} है और $Q(0),
Q(1), \dots, Q(n) \in \Z$। $(B_k)_{k \leq n}$ में उसके \omterm{prop:b1:vspaces:coordinates}{निर्देशांक} $c_k =
\sum_{j \leq k}(-1)^{k-j}\binom kj Q(j) \in \Z$ हैं (प्रश्न 14 केवल $0,
\dots, k \leq n$ पर के मानों का प्रयोग करता है)। प्रश्न 17 ($\Leftarrow$) से
$Q$ पूरे $\Z$ पर पूर्णांक-मान है, अतः $P(X) = Q(X - a)$ भी।

\textbf{19.} $k$ क्रमागत पूर्णांकों का गुणनफल किसी $m \in \Z$ के लिए $m(m-1)
\cdots(m-k+1) = k!\,B_k(m)$ है, और प्रश्न 16 से $B_k(m) \in \Z$: अतः गुणनफल
$k!$ से विभाज्य है।

\textbf{20.} $0,1,2,3$ पर $P = \frac{X(X+1)(2X+1)}{6}$ के मान: $0, 1, 5,
14$। अंतर सारणी: $\Delta$ की पंक्ति $1, 4, 9$; $\Delta^2$ की पंक्ति $3, 5$;
$\Delta^3$ की पंक्ति $2$। अतः प्रश्न 13 से,
\[
P = 0\cdot B_0 + 1\cdot B_1 + 3\,B_2 + 2\,B_3 ,
\]
और \omterm{prop:b1:vspaces:coordinates}{निर्देशांक} पूर्णांक हैं: $P$ पूर्णांक-मान है (प्रश्न 17), जबकि उसके एकपदी
गुणांक $\frac13, \frac12, \frac16$ पूर्णांक नहीं हैं। सीधी संगणना:
\begin{align*}
\Delta P &= \frac{(X+1)(X+2)(2X+3) - X(X+1)(2X+1)}{6} \\
&= \frac{(X+1)\bigl[(X+2)(2X+3) - X(2X+1)\bigr]}{6}
= \frac{(X+1)(6X+6)}{6} = (X+1)^2 .
\end{align*}
$P(m) = \sum_{j=0}^{m-1}\Delta P(j) = \sum_{j=1}^{m} j^2$ को दूरबीनी करने पर
($P(0) = 0$ के साथ): अर्थात् वर्गों के योग का सूत्र।

\textbf{21.} $i = 0, \dots, n$ पर मान $2^i$ की अंतर सारणी बाएँ किनारे पर
लगातार $1$ है: अनुक्रम $(2^i)$ का $\Delta^k$ फिर से $(2^i)$ है (क्योंकि
$2^{i+1} - 2^i = 2^i$), अतः सभी $k \leq n$ के लिए $\bigl(\Delta^k P\bigr)(0)
= 2^0 = 1$, और प्रश्न 13 से $P = B_0 + B_1 + \dots + B_n$। तब
\[
P(n+1) = \sum_{k=0}^{n}\binom{n+1}{k}
= 2^{n+1} - \binom{n+1}{n+1} = 2^{n+1} - 1 \neq 2^{n+1}:
\]
अर्थात् प्रतिरूप पहले ही अनियंत्रित बिंदु पर टूट जाता है।

\textbf{22.} प्रश्न 12 से $B_k = \Delta B_{k+1}$, अतः
\[
\sum_{j=0}^{m-1} B_k(j)
= \sum_{j=0}^{m-1}\bigl(B_{k+1}(j+1) - B_{k+1}(j)\bigr)
= B_{k+1}(m) - B_{k+1}(0) = B_{k+1}(m).
\]
$j < k$ के लिए पद $B_k(j)$ शून्य हो जाते हैं, अतः योग वस्तुतः $j = k$ से
आरंभ होता है: $\sum_{j=k}^{m-1}\binom jk = \binom m{k+1}$, अर्थात्
हॉकी-स्टिक सर्वसमिका।

\textbf{23.} अंतर सारणियाँ (अथवा सीधा प्रसार) देती हैं
\[
X^2 = B_1 + 2 B_2, \qquad X^3 = B_1 + 6 B_2 + 6 B_3
\]
(जाँच: $B_1 + 2B_2 = X + X(X-1) = X^2$; $X = 1, 2, 3$ पर दूसरा $1, 8, 27$
देता है)। फिर प्रश्न 22 से मिलता है
\[
\sum_{j=0}^{m-1} j^2 = B_2(m) + 2B_3(m)
= \binom m2 + 2\binom m3 = \frac{m(m-1)(2m-1)}{6},
\]
\[
\sum_{j=0}^{m-1} j^3 = B_2(m) + 6B_3(m) + 6B_4(m)
= \binom m2 + 6\binom m3 + 6\binom m4 .
\]
अंतिम व्यंजक को प्रसारित करने पर: $\binom m2 + 6\binom m3 + 6\binom m4 =
\frac{m(m-1)}{2}\bigl[1 + 2(m-2) + \frac{(m-2)(m-3)}{2}\bigr] =
\frac{m^2(m-1)^2}{4} = \binom m2^2$। $m$ के स्थान पर $m + 1$ रखने पर: $1^3 +
\dots + m^3 = \bigl(\frac{m(m+1)}2\bigr)^2 = (1 + \dots + m)^2$, अर्थात्
निकोमैकस की सर्वसमिका।

\textbf{24.} गाँठें $0, 1, 2$, \omterm{def:b1:poly:def}{बहुपद} $X^2$। एकपदी \omterm{def:b1:vspaces:free}{आधार} $(1, X, X^2)$:
\omterm{prop:b1:vspaces:coordinates}{निर्देशांक} $(0, 0, 1)$। लाग्रांज \omterm{def:b1:vspaces:free}{आधार}: \omterm{prop:b1:vspaces:coordinates}{निर्देशांक} मान $(0, 1, 4)$ हैं
(प्रश्न 4)। न्यूटन \omterm{def:b1:vspaces:free}{आधार} $(1, X, X(X-1))$: \omterm{pb:b1:vspaces:1}{विभाजित अंतर} $f[0] = 0$, $f[0,1] =
1$, $f[0,1,2] = \frac{3 - 1}{2} = 1$ (प्रश्न 9), अतः \omterm{prop:b1:vspaces:coordinates}{निर्देशांक} $(0, 1, 1)$
--- और वस्तुतः $X + X(X-1) = X^2$। तीन \omterm{def:b1:vspaces:free}{आधार}, तीन निर्देशांक-सदिश, एक ही
\omterm{def:b1:poly:def}{बहुपद}।

\textbf{25.} (क) प्रश्न 4 स्वतः इसलिए है कि $(L_i)$ एक \emph{\omterm{def:b1:vspaces:free}{आधार}} है:
अंतर्वेशन का अस्तित्व और अद्वितीयता ठीक निर्देशांकों का अस्तित्व और
अद्वितीयता ही हैं। (ख) न्यूटन \omterm{def:b1:vspaces:free}{आधार} एक सीढ़ी है, अतः \omterm{prop:b1:vspaces:coordinates}{निर्देशांक} क्रमागत भागों
से संगणित होते हैं --- हर नई गाँठ पिछली गाँठों को छेड़े बिना एक पद जोड़ देती
है --- जबकि $P$ के लाग्रांज \omterm{prop:b1:vspaces:coordinates}{निर्देशांक} तो मान $P(x_i)$ ही हैं, जो बिना किसी
संगणना के उपलब्ध हैं। (ग) दोनों \omterm{def:b1:vspaces:free}{आधार} \cref{prop:b1:vspaces:freecriteria} की
एक ही जोड़ी कसौटियों से स्वतंत्र हैं: न्यूटन के लिए भिन्न-भिन्न घातें, और
लाग्रांज के लिए गाँठों पर मूल्यांकन। (घ) पोल्या की प्रमेय कहती है कि
``$P(\Z) \subseteq \Z$'', जो मानों का गुण है, \omterm{def:b1:vspaces:free}{आधार} $(B_k)$ में निर्देशांकों
की पूर्णांकता के तुल्य है --- किसी \omterm{def:b1:poly:def}{बहुपद} का अंकगणित केवल उसी \omterm{def:b1:vspaces:free}{आधार} में दिखाई
देता है जो प्रश्न के अनुकूल हो।
\end{solution}
